\section{Yorug‘lik  difraksiya}


\begin{enumerate}

\item Difraksiya deb qanday hodisaga aytiladi?
\begin{enumerate}[label=\Alph*)]
\item o‘lchamlari to‘lqin uzunliklari tartibida bo‘lgan to‘siq yoki tirqishga tushgan to‘lqinlarning to‘silgan sohaga og‘ib tarqalish hodisasi.
\item ikki kogerent to‘lqin qo‘shilganda hosil bo‘lgan to‘lqin amplitudasining kuchayishi yoki susayishi hodisasi.
\item to‘siqqa uchragan to‘lqinlarning to‘siqdan biror burchak ostida orqaga qaytib tarqalish hodisasi.
\item ikki to‘lqinning qo‘shilishi natijasida natijaviy tebranishlar amplitudasining keskin ortib yoki kamayib ketish hodisasi.
\item TJY.
\end{enumerate}

\item Yorug‘lik to‘lqinlarining tarqalishda to‘g‘ri chiziqdan og‘ishi va to‘siqni aylanib, soya sohasiga o‘tishi ... deb ataladi.
\begin{enumerate}[label=\Alph*)]
\item interferensiya.
\item yorug‘likning qutblanishi.
\item fotoeffekt.
\item dispersiya.
\item difraksiya.
\end{enumerate}

\item Yorug‘lik qanday sharoitda to‘g‘ri chiziq bo‘ylab tarqalishdan chetlashadi?
\begin{enumerate}[label=\Alph*)]
\item bir jinsli bo‘lmagan muhitda.
\item o‘lchamlari yorug‘lik to‘lqin uzunligi bilan solishtirarli to‘siqlarga uchraganda.
\item o‘lchamlari yorug‘lik to‘lqin uzunligi bilan solishtirarli tirqish yoki teshiklardan o‘tganda.
\item TJY.
\item A, B va C javoblar to‘g‘ri.
\end{enumerate}

\item Uzunligi 800 nm bo‘lgan ikki kki kogerent yorug‘lik to‘lqini fazoning bir nuqtasiga 2 $\mu$m optik yo‘l farqi bilan kelib tushmoqda. Nurlar kelib tushayotgan nuqtada interferensiya hodisasi bo‘ladimi yoki difraksiyami? Shu hodisa tufayli to‘lqin bu nuqtada kuchayadimi yoki susayadimi?
\begin{enumerate}[label=\Alph*)]
\item interferensiya, susayadi.
\item difraksiya, susayadi.
\item interferensiya, kuchayadi.
\item difraksiya, susayadi.
\item TJY.
\end{enumerate}

\item Fazoning biror-bir nuqtasiga to‘lqin uzunligi 600 nm, optik yo‘llar farqi 1,2 $\mu$m bo‘lgan kogerent nurlar kelib tushmoqda. Bu nuqtada qanday hodisa kuzatiladi?
\begin{enumerate}[label=\Alph*)]
\item difraksion maksimumlar.
\item interferension minimumlar.
\item interferension maksimumlar.
\item difraksion minimumlar.
\item TJY.
\end{enumerate}

\item Agar difraksion panjaraning bitta shtrixining kengligi $a=1{,}25$ $\mu$m va shtrixlar orasidagi masofa $b=2{,}35$ $\mu$m bo‘lsa, panjara davri (doimiysi) $d$ nimaga teng($\mu$m)?

\begin{tikzpicture}[>=Stealth, line cap=round, font=\small]

% Difraksion panjara shtrixlari
\foreach \x in {0,1.15,2.30,3.45,4.60}{
    \fill[gray!50] (\x,0) rectangle ++(0.48,1.75);
    \draw[thick] (\x,0) -- (\x,1.75);
    \draw[thick] (\x+0.48,0) -- (\x+0.48,1.75);
}

% Pastki tayanch chizig'i
\draw[very thick] (0,0) -- (5.08,0);

% Belgilash uchun yordamchi punktir chiziqlar
\draw[densely dashed] (1.15,-0.08) -- (1.15,-0.68);
\draw[densely dashed] (1.63,-0.08) -- (1.63,-0.68);
\draw[densely dashed] (2.30,-0.08) -- (2.30,-0.68);

% a kenglik
\draw[<->] (1.15,-0.68) -- node[below=2pt] {$a$} (1.63,-0.68);

% b oraliq
\draw[<->] (1.63,-0.68) -- node[below=2pt] {$b$} (2.30,-0.68);

\end{tikzpicture}

\begin{enumerate}[label=\Alph*)]
\item 1,1.
\item 1,88.
\item 0,53.
\item 3,6.
\item 1,71.
\end{enumerate}

\item 1 mm da 1000 ta shtrixi bor difraksion panjara davriysini ($\mu$m) aniqlang.
\begin{enumerate}[label=\Alph*)]
\item 100.
\item 10.
\item 2.
\item 0,1.
\item 1.
\end{enumerate}

\item Zamonaviy difraksion panjarada 3/4 mm masofada 12500 ta shtrix chizilgan. Panjara doimiysi ($\mu$m) qanday?
\begin{enumerate}[label=\Alph*)]
\item 0,03.
\item 0,24.
\item 0,12.
\item 0,06.
\item 0,1.
\end{enumerate}

\item Agar panjara davri 2 $\mu$m bo‘lsa, u holda har bir millimetrda necha shtrix mavjud?
\begin{enumerate}[label=\Alph*)]
\item 500.
\item 400.
\item 300.
\item 200.
\item 100.
\end{enumerate}

\item Difraksion panjarada ikkita qo‘shni tirqishlarning chetlaridan difraksiyaga uchragan to‘lqinlar yo‘l farqi 2,5 $\mu$m va difraksiya burchagi $30^\circ$ bo‘lsa, shu panjaraning doimiysi nech $\mu$m ga teng?

\begin{tikzpicture}[>=Stealth, line cap=round, font=\small]

% Manba
\coordinate (S) at (0,1.28);

\foreach \a in {0,45,...,315}
    \draw[red] ($(S)+(\a:0.13)$) -- ($(S)+(\a:0.27)$);
\fill[red] (S) circle (0.105);

\node[above=8pt] at (S) {Manba};

% Difraksion panjara: ikkita qo'shni tirqish
\fill (2.05,1.88) rectangle (2.18,2.62);
\fill (2.05,0.92) rectangle (2.18,1.60);
\fill (2.05,0.00) rectangle (2.18,0.65);

% Tirqishlar orasidagi masofa d
\draw[<->,thick] (1.78,0.77) -- node[left=5pt] {$d$} (1.78,1.74);

% Ekran
\draw[very thick] (5.80,0) -- (5.80,2.68);
\fill (5.80,1.28) circle (2pt);
\fill (5.80,2.68) circle (2pt);
\node[below=6pt] at (5.80,0) {Ekran};

% Markaziy nur
\draw (2.18,1.28) -- (5.80,1.28);

% Difraksiyalangan nurlar
\draw[densely dashed,red] (2.18,1.74) -- (5.80,2.68);
\draw[densely dashed,red] (2.18,0.77) -- (5.80,2.68);

% Difraksiya burchagi
\draw (3.80,1.28) arc[start angle=0,end angle=30,radius=0.58];
\node at (4.1,1.51) {$30^\circ$};

% Yo'llar farqi belgisi
\coordinate (A) at (2.18,0.77);
\coordinate (B) at (2.6,1);
\draw[thick] (A) -- (B);
\draw (2.18,1.74)--(B);
\draw[<->] (2.18,0.6)--(2.65,0.85);



\node[below right] at (2.20,0.7) {$\Delta l=2{,}5\ \mu\mathrm{m}$};

\end{tikzpicture}

\begin{enumerate}[label=\Alph*)]
\item 1,25.
\item 2,5.
\item 5.
\item 10.
\item 25.
\end{enumerate}

\item Davri 1,5 $\mu$m bo‘lgan difraksion panjara $30^\circ$ burchakka og‘dirgan birinchi tartibli spektr chiqishga mos keladigan to‘lqin uzunligini toping (nm).
\begin{enumerate}[label=\Alph*)]
\item 300.
\item 450.
\item 500.
\item 600.
\item 750.
\end{enumerate}

\item To‘lqin uzunligi 600 nm bo‘lgan monoxromatik yorug‘lik to‘lqini davri 3,6 $\mu$m bo‘lgan difraksion panjaraga tik tushmoqda. $30^\circ$ burchak ostida ko‘ringan difraksion maksimumning tartibini aniqlang.
\begin{enumerate}[label=\Alph*)]
\item 1.
\item 2.
\item 3.
\item 4.
\item 5.
\end{enumerate}

\item Davriysi $d$ bo‘lgan difraksion panjaraga to‘lqin uzunligi $\lambda$ bo‘lgan monoxromatik parallel nurlar dastasi tik tushmoqda. Birinchi bosh maksimum burchagi $\varphi$ uchun qaysi shart bajarilishi kerak?
\begin{enumerate}[label=\Alph*)]
\item $\sin\varphi=\lambda/d$.
\item $\sin\varphi=d/\lambda$.
\item $\cos\varphi=\lambda/d$.
\item $\cos\varphi=d/\lambda$.
\item $\cos\varphi=\lambda d$.
\end{enumerate}

\item Davri 0,001 mm bo‘lgan difraksion panjara yordamida hosil qilingan birinchi tartibli spektrda yashil $(\lambda=0{,}5~\mu\text{m})$ chiziq qanday burchak ostida ko‘rinadi?
\begin{enumerate}[label=\Alph*)]
\item $0^\circ$.
\item $30^\circ$.
\item $45^\circ$.
\item $60^\circ$.
\item $75^\circ$.
\end{enumerate}

\item Difraksion panjara doimiysi 1,4 $\mu$m. Panjarada hosil qilingan difraksion manzaraning 1-tartibli maksimumi $30^\circ$ burchak ostida kuzatilayotgan bo‘lsa, nurning to‘lqin uzunligi ($\mu$m) ni toping.
\begin{enumerate}[label=\Alph*)]
\item 0,5.
\item 0,7.
\item 0,6.
\item 0,8.
\item 1.
\end{enumerate}

\item Davri $1\cdot10^{-3}$ mm bo‘lgan difraksion panjarada kuzatilayotgan birinchi tartibli simmetrik maksimumlar orasida burchak $60^\circ$ bo‘lsa, tushayotgan yorug‘likning to‘lqin uzunligi qanday (nm) bo‘ladi?
\begin{enumerate}[label=\Alph*)]
\item 400.
\item 500.
\item 600.
\item 700.
\item 870.
\end{enumerate}

\item Difraksion panjara ortidagi 6-tartibli maksimum $60^\circ$ burchak ostida kuzatiladi. Panjara doimiysining nurning to‘lqin uzunligiga nisbatini toping.
\begin{enumerate}[label=\Alph*)]
\item 9.
\item 7,2.
\item 5,2.
\item 6,9.
\item 3,6.
\end{enumerate}

\item Davri 0,0025 mm bo‘lgan difraksion panjara yordamida hosil qilingan 3-tartibli spektrdagi zangori $(\lambda=0{,}5~\mu\text{m})$ rangli nurlarning og‘ish burchagi qanday?
\begin{enumerate}[label=\Alph*)]
\item $37^\circ$.
\item $45^\circ$.
\item $90^\circ$.
\item $60^\circ$.
\item $30^\circ$.
\end{enumerate}

\item To‘lqin uzunligi 500 nm bo‘lgan monoxromatik yorug‘lik to‘lqini $\sqrt3~\mu$m bo‘lgan difraksion panjaraga tik tushmoqda. Uchinchi tartibli difraksion maksimum qanday burchak ostida ko‘rinadi?
\begin{enumerate}[label=\Alph*)]
\item $90^\circ$.
\item $60^\circ$.
\item $45^\circ$.
\item $30^\circ$.
\item $20^\circ$.
\end{enumerate}









\item Difraksion panjaraning har 1 mm ida 250 ta shtrix bor. Difraksion panjaraga to‘lqin uzunligi 500 nm bo‘lgan monoxromatik yorug‘lik tushadi. To‘rtinchi difraksion maksimum qanday burchak ostida kuzatiladi?
\begin{enumerate}[label=\Alph*)]
\item $75^\circ$.
\item $60^\circ$.
\item $45^\circ$.
\item $30^\circ$.
\item $15^\circ$.
\end{enumerate}

\item To‘lqin uzunligi $\lambda$ bo‘lgan yorug‘likning davri $5\lambda$ bo‘lgan difraksion panjarasida difraksiyalanishida 2-tartibli maksimum qanday burchak ostida kuzatiladi?
\begin{enumerate}[label=\Alph*)]
\item $\arcsin(3/5)$.
\item $\arcsin(0{,}45)$.
\item $\arcsin(0{,}3)$.
\item $\arcsin(0{,}15)$.
\item $\arcsin(2/5)$.
\end{enumerate}

\item Difraksion panjaraning 1 mm masofasida 500 ta shtrix bo‘lsa va ikkinchi tartibli maksimum $30^\circ$ burchak ostida kuzatilayotgan bo‘lsa, yorug‘likning to‘lqin uzunligi qanday ($\mu$m)?
\begin{enumerate}[label=\Alph*)]
\item 0,5.
\item 0,6.
\item 0,7.
\item 0,8.
\item 0,9.
\end{enumerate}

\item Monoxromatik yorug‘lik dastasi difraksion panjaraga tik tushadi. Uchinchi tartibli maksimum markazga nisbatan $30^\circ$ burchak ostida kuzatiladi. Panjara davrini $d$ ni yorug‘likning to‘lqin uzunligi $\lambda$ bilan ifodalang.
\begin{enumerate}[label=\Alph*)]
\item $d=\lambda$.
\item $d=1{,}5\lambda$.
\item $d=3\lambda$.
\item $d=6\lambda$.
\item $d=\lambda/3$.
\end{enumerate}

\item To‘lqin uzunligi $\lambda$ bo‘lgan monoxromatik parallel yorug‘lik nurlari davri $d=6\lambda$ bo‘lgan difraksion panjaraga normal ravishda tushadi. Ikkinchi tartibli maksimum kuzatiladigan difraksiya burchagi $\alpha$ bo‘lsa, $\sin\alpha$ nimaga teng?

\begin{tikzpicture}[
    >=Stealth,
    line cap=round,
    font=\small
]

% Parallel tushuvchi nurlar
\foreach \y in {0.62,0.82,1.02,1.22}{
    \draw[->,thick,red] (-1.55,\y) -- (-0.55,\y);
}
\node[above] at (-1.05,1.28) {$\lambda$};

% Difraksion panjara
\fill (-0.10,1.62) rectangle (0.02,2.25);
\fill (-0.10,0.92) rectangle (0.02,1.43);
\fill (-0.10,0.18) rectangle (0.02,0.72);
\fill (-0.10,-0.55) rectangle (0.02,-0.02);

\node[above, align=center] at (-0.04,2.30) {difraksion\\panjara};
\node[below] at (-0.04,-0.58) {$d=6\lambda$};

% Markaziy va og'gan nurlar
\coordinate (O) at (0.03,0.82);
\coordinate (Kzero) at (3.10,0.82);
\coordinate (Kone) at (3.10,1.52);
\coordinate (Ktwo) at (3.10,2.22);

\draw[thick] (O) -- (Kzero);
\draw[densely dashed] (O) -- (Kone);
\draw[densely dashed] (O) -- (Ktwo);

% Alfa burchagi
\draw (0.72,0.82) arc[start angle=0,end angle=25,radius=0.69];
\node at (0.92,1.02) {$\alpha$};

% Ekran: maksimumlar
\draw[thick] (3.10,0.18) rectangle (3.30,2.38);

\fill[red] (3.10,0.63) rectangle (3.30,1.01);
\fill[red] (3.10,1.34) rectangle (3.30,1.70);
\fill[red] (3.10,2.05) rectangle (3.30,2.37);

% Ekrandagi bo'linishlar
\draw (3.10,0.63) -- (3.30,0.63);
\draw (3.10,1.01) -- (3.30,1.01);
\draw (3.10,1.34) -- (3.30,1.34);
\draw (3.10,1.70) -- (3.30,1.70);
\draw (3.10,2.05) -- (3.30,2.05);

% Tartib raqamlari
\node[right=5pt] at (3.30,0.82) {$k=0$};
\node[right=5pt] at (3.30,1.52) {$k=1$};
\node[right=5pt] at (3.30,2.22) {$k=2$};

% Ekran yozuvi
\node[below=5pt] at (3.20,0.18) {Ekran};

\end{tikzpicture}

\begin{enumerate}[label=\Alph*)]
\item $1/4$.
\item $1/3$.
\item $2{,}3$.
\item $1/6$.
\item 0.
\end{enumerate}

\item Bir 2,5 mm ida 4000 ta shtrixi bo‘lgan difraksion panjaraga to‘lqin uzunligi 500 nm bo‘lgan monoxromatik yorug‘lik tushmoqda. Birinchi tartibli ikki spektr orasidagi burchakni toping.
\begin{enumerate}[label=\Alph*)]
\item $60^\circ$.
\item $90^\circ$.
\item $30^\circ$.
\item $106^\circ$.
\item $45^\circ$.
\end{enumerate}

\item Natriy sariq chizig‘ining to‘lqin uzunligi 500 nm. Difraksion panjaraning davri 2,8 $\mu$m bo‘lsa, spektrning eng katta tartibi qanday?
\begin{enumerate}[label=\Alph*)]
\item 6.
\item 4.
\item 5.
\item 3.
\item 2.
\end{enumerate}

\item Difraksiya panjarasining davri 2 $\mu$m. Natriyning to‘lqin uzunligi 589 nm sariq chizig‘ining difraksiya manzarasidagi eng katta tartib nomerini toping.
\begin{enumerate}[label=\Alph*)]
\item 6.
\item 5.
\item 4.
\item 3.
\item 2.
\end{enumerate}

\item Har bir millimetrida 200 ta shtrixi bo‘lgan difraksion panjaraga to‘lqin uzunligi $5\cdot10^{-5}$ cm bo‘lgan monoxromatik to‘lqin normal tushmoqda. Kuzatish mumkin bo‘lgan difraksion maksimumlarning eng katta tartibi qanday?
\begin{enumerate}[label=\Alph*)]
\item 1.
\item 5.
\item 10.
\item 20.
\item 100.
\end{enumerate}

\item 1 mm ga 400 shtrix to‘g‘ri keladigan difraksion panjaraga to‘lqin uzunligi $5\cdot10^{-5}$ cm bo‘lgan monoxromatik yorug‘lik to‘lqini tik tushmoqda. Spektrning eng katta tartibini aniqlang.
\begin{enumerate}[label=\Alph*)]
\item 3.
\item 7.
\item 5.
\item 6.
\item 2.
\end{enumerate}

\item To‘lqin uzunligi 0,5 $\mu$m bo‘lgan monoxromatik yorug‘lik 2,9 $\mu$m bo‘lgan difraksion panjaraga tik tushmoqda. Ekranda nechta bosh maksimum kuzatiladi?
\begin{enumerate}[label=\Alph*)]
\item 5.
\item 6.
\item 10.
\item 11.
\item 12.
\end{enumerate}

\item Doimiysi 1,1 $\mu$m bo‘lgan yassi difraksion panjaraga to‘lqin uzunligi 0,5 $\mu$m bo‘lgan monoxromatik to‘lqin normal tushmoqda. Kuzatish mumkin bo‘lgan maksimumlar sonini toping.
\begin{enumerate}[label=\Alph*)]
\item 9.
\item 7.
\item 5.
\item 4.
\item 2.
\end{enumerate}

\item Davri 3,4 $\mu$m bo‘lgan difraksion panjaraga tik ravishda to‘lqin uzunligi 0,5 $\mu$m bo‘lgan monoxromatik yorug‘lik tushmoqda. Ekranda nechta bosh maksimum kuzatiladi?
\begin{enumerate}[label=\Alph*)]
\item 6.
\item 7.
\item 12.
\item 14.
\item 13.
\end{enumerate}

\item Difraksion panjara davri $d=2\lambda$ bo‘lsa, panjara ortida nechta difraksion maksimum kuzatiladi?
\begin{enumerate}[label=\Alph*)]
\item 2.
\item 1.
\item 5.
\item 3.
\item 4.
\end{enumerate}

\item Yashil $(\lambda=0{,}5~\mu\text{m})$ nur bilan yoritilgan davriysi 100 $\mu$m bo‘lgan difraksion panjara nechta difraksion maksimum hosil qiladi?
\begin{enumerate}[label=\Alph*)]
\item 500.
\item 401.
\item 200.
\item 50.
\item 10.
\end{enumerate}

\item To‘lqin uzunligi 0,5 $\mu$m bo‘lgan yorug‘lik dastasi doimiysi 2,3 $\mu$m bo‘lgan difraksion panjaraga perpendikulyar tushadi. Ekranda kuzatiladigan maksimumlarning maksimal soni ...
\begin{enumerate}[label=\Alph*)]
\item 9.
\item 5.
\item 4.
\item 2.
\item 1.
\end{enumerate}

\item Monoxromatik yorug‘lik dastasi difraksion panjaraga tik tushadi. To‘rtinchi tartibli maksimum $30^\circ$ ostida kuzatiladi. Panjara tomonidan berilgan maksimumlarning umumiy soni ...
\begin{enumerate}[label=\Alph*)]
\item 7.
\item 8.
\item 15.
\item 17.
\item 19.
\end{enumerate}

\item Difraksion panjaradan ekrangacha bo‘lgan masofa 80 cm, nolinchi maksimumdan uchinchi maksimumgacha masofa 3,6 mm bo‘lsa, yorug‘likning to‘lqin uzunligi necha metrga teng? Panjara doimiysi 0,3 mm ga teng.
\begin{enumerate}[label=\Alph*)]
\item $4{,}5\cdot10^{-7}$.
\item $8{,}7\cdot10^{-7}$.
\item $6\cdot10^{-7}$.
\item $0{,}65\cdot10^{-7}$.
\item $8\cdot10^{-7}$.
\end{enumerate}

\item Davri 0,02 mm bo‘lgan difraksion panjara yordamida markaziy maksimumdan 3 cm masofada birinchi maksimum hosil qilindi. Panjaradan ekrangacha bo‘lgan masofa 1 m bo‘lsa, yorug‘likning to‘lqin uzunligi qanday ($\mu$m)?
\begin{enumerate}[label=\Alph*)]
\item 0,4.
\item 0,8.
\item 0,86.
\item 0,2.
\item 0,6.
\end{enumerate}

\item 1 mm da 100 ta shtrixi bo‘lgan difraksion panjaraga yorug‘lik tushmoqda. Ekrandagi birinchi maksimum nolinchi maksimumdan 12 cm uzoqlikda kuzatiladi. Ekran bilan difraksion panjara orasidagi masofa 2 m bo‘lsa, yorug‘likning to‘lqin uzunligi qanday ($\mu$m)?
\begin{enumerate}[label=\Alph*)]
\item 60.
\item 6,5.
\item 6,4.
\item 0,7.
\item 0,6.
\end{enumerate}

\item Difraksion panjara vositasida 1,8 m masofada, markaziy maksimumdan 5,4 cm masofada birinchi tartibli maksimum hosil qilinmoqda. Nurning to‘lqin uzunligi 0,6 $\mu$m bo‘lsa, panjara doimiysi ($\mu$m) qanday?
\begin{enumerate}[label=\Alph*)]
\item 15.
\item 40.
\item 20.
\item 30.
\item 10.
\end{enumerate}


\item Ekrandan panjaragacha bo‘lgan $a$ masofa bilan nolinchi maksimumdan birinchi maksimumgacha bo‘lgan $b$ masofa orasida qanday bog‘lanish bor?
\begin{enumerate}[label=\Alph*)]
\item $a$ ortsa $b$ o‘zgarmaydi.
\item $a$ ortsa $b$ kamayadi.
\item $a$ ortsa $b$ ham ortadi.
\item $a$ va $b$ bir-biriga bog‘liq emas.
\item TJY.
\end{enumerate}

\item Difraksiya maksimumi bilan ekran markazi orasidagi $b$ masofaning panjara davri $d$ ga bog‘lanishi qanday?
\begin{enumerate}[label=\Alph*)]
\item $a$ ortsa, $b$ kamayadi.
\item $a$ ortsa, $b$ ortadi.
\item $a$ ortsa, $b$ o‘zgarmaydi.
\item $b$ $a$ ga bog‘liq emas.
\item TJY.
\end{enumerate}

\item Davri kichik bo‘lgan difraksion panjara o‘rniga davri kattaroq bo‘lgan panjara olinsa, ekrandagi maksimumlar orasidagi masofa qanday o‘zgaradi? Maksimumlar soni-chi?
\begin{enumerate}[label=\Alph*)]
\item qisqaradi, kamayadi.
\item uzayadi, ko‘payadi.
\item o‘zgarmaydi, kamayadi.
\item uzayadi, kamayadi.
\item qisqaradi, ko‘payadi.
\end{enumerate}

\item Har 1 millimetrida $N=25$ ta shtrixi bo‘lgan difraksion panjaraga monoxromatik yorug‘lik tik tushmoqda. Ekrandagi $k=4$ tartibli ikkita bosh maksimum orasidagi masofa 5,4 cm, ekran bilan difraksion panjara oralig‘i 50 cm bo‘lsa, yorug‘likning to‘lqin uzunligi qanday ($\mu$m)?
\begin{enumerate}[label=\Alph*)]
\item 0,52.
\item 0,54.
\item 0,56.
\item 0,58.
\item 0,5.
\end{enumerate}

\item Agar monoxromatik yorug‘lik davri $d$ bo‘lgan difraksiya panjarasida difraksiyalanganda 1- tartibli maksimum ekranda markaziy maksimumdan $x$ masofada joylashgan bo‘lsa, yorug‘likning to‘lqin uzunligi qanday? Ekran va panjara orasidagi masofa $L$.


\begin{enumerate}[label=\Alph*)]
\item $\dfrac{dx}{\sqrt{L^2+x^2}}$.
\item $\dfrac{dL}{\sqrt{L^2+x^2}}$.
\item $\dfrac{dx}{L}$.
\item $\dfrac{d\sqrt{L^2+x^2}}{L}$.
\item $\dfrac{d\sqrt{L^2+x^2}}{x}$.
\end{enumerate}

\item Difraksion panjarasi yordamida yorug‘lik difraksiyalanganda 3-tartibli spektrdagi $\lambda_1=780$ nm to‘lqin uzunlikli chiziq bilan 4-tartibli spektrdagi qanday $\lambda_2$ to‘lqin uzunlikli chiziq ustma-ust tushadi?

\begin{tikzpicture}[
    >=Stealth,
    line cap=round,
    font=\small
]

% Parallel tushuvchi nurlar
\foreach \y in {0.62,0.82,1.02,1.22}{
    \draw[->,thick,red] (-1.55,\y) -- (-0.55,\y);
}
\node[above] at (-1.1,1.28) {Oq yoru`lik};

% Difraksion panjara
\fill (-0.10,1.62) rectangle (0.02,2.25);
\fill (-0.10,0.92) rectangle (0.02,1.43);
\fill (-0.10,0.18) rectangle (0.02,0.72);
\fill (-0.10,-0.55) rectangle (0.02,-0.02);

\node[above, align=center] at (-0.04,2.30) {difraksion\\panjara};
\node[below] at (-0.04,-0.58) {$d$};

% Markaziy va og'gan nurlar
\coordinate (O) at (0.03,0.82);
\coordinate (Kzero) at (4.10,0.82);
\coordinate (Kfour) at (4.10,3);
\coordinate (Ktwo) at (3.10,2.22);
\coordinate (Kthree) at (3.10,2.92);

\draw[thick] (O) -- (Kzero);
\draw[densely dashed] (O) -- (Kfour);
\draw[densely dashed] (O) -- (Kthree);
\draw[densely dashed] (O) -- (Kthree);

% Alfa burchagi
\draw (0.72,0.82) arc[start angle=0,end angle=25,radius=0.69];
\node at (0.92,1.02) {$\alpha$};

% Ekran: maksimumlar
\draw[thick] (3.10,0.18) rectangle (3.30,3.38);

\fill[red] (3.10,0.63) rectangle (3.30,1.01);
\fill[red] (3.10,1.34) rectangle (3.30,1.70);
\fill[red] (3.10,2.05) rectangle (3.30,2.37);
\fill[red] (3.10,2.72) rectangle (3.30,3.04);


% Ekran: maksimumlar
\draw[thick] (4.10,0.18) rectangle (4.30,3.38);

\fill[red] (4.10,0.72) rectangle (4.30,1);
\fill[red] (4.10,1.2) rectangle (4.30,1.42);
\fill[red] (4.10,1.65) rectangle (4.30,1.9);
\fill[red] (4.10,2.2) rectangle (4.30,2.5);
\fill[red] (4.10,2.8) rectangle (4.30,3.1);



% Ekrandagi bo'linishlar
\draw (3.10,0.63) -- (3.30,0.63);
\draw (3.10,1.01) -- (3.30,1.01);
\draw (3.10,1.34) -- (3.30,1.34);
\draw (3.10,1.70) -- (3.30,1.70);
\draw (3.10,2.05) -- (3.30,2.05);

% Tartib raqamlari

\node[right=2pt] at (2,1.52) {$k=1$};
\node[right=2pt] at (2,2.22) {$k=2$};
\node[right=2pt] at (2,2.92) {$k=3$};


\node[right=6pt] at (4,1.35) {$k=1$};
\node[right=6pt] at (4,1.82) {$k=2$};
\node[right=6pt] at (4,2.38) {$k=3$};
\node[right=6pt] at (4,2.94) {$k=4$};

% Ekran yozuvi
\node[below=5pt] at (3.70,0.18) {Ekran};

\end{tikzpicture}




\begin{enumerate}[label=\Alph*)]
\item 292.
\item 347.
\item 520.
\item 585.
\item 1040.
\end{enumerate}

\item Qizil va binafsha nurlar difraksion panjarasiga tushadi. Quyidagi mulohazalardan qaysi to‘g‘ri?
\begin{enumerate}
\item[1] qizil nurning 3-maksimumi binafsha nurning 3-maksimumiga qaraganda markaziy maksimumdan uzoqroqda joylashgan.
\item[2] qizil va binafsha ranglarning markaziy maksimumlari bir nuqtada joylashadi.
\item[3] qizil va binafsha ranglarning markaziy maksimumlari bir nuqtada joylashmaydi.
\item[4] binafsha nurning 3-maksimumi qizil nurning 3-maksimumiga qaraganda markaziy maksimumdan uzoqroqda joylashgan.
\item[5] qizil nurning jami maksimumlari soni binafsha nurning maksimumlari sonidan ko‘p bo‘ladi.
\end{enumerate}
\begin{enumerate}[label=\Alph*)]
\item 1 va 3.
\item 3 va 4.
\item 3 va 5.
\item 1, 2 va 4.
\item 1, 2 va 5.
\end{enumerate}

\item Maktabda 1 mm da 50 va 100 ta shtrixi bo‘lgan difraksion panjara bor. Bir xil sharoitda nurlarning qaysi biri ekranda enlqroq spektr hosil qiladi?
\begin{enumerate}[label=\Alph*)]
\item spektrlari bir xil bo‘ladi.
\item spektrning eni shtrixlar soniga bog‘liq emas.
\item birinchisi.
\item ikkinchisi.
\item TJY.
\end{enumerate}

\item Nomlari quyida keltirilgan asboblarning qaysilari oq yorug‘likni tarkibiy qismlarga (spektrlarga) ajratish uchun xizmat qiladi?
\begin{enumerate}
\item yassi ko‘zgu;
\item yig‘uvchi yupqa linza;
\item difraksion panjara;
\item shaffof uch qoli prizma.
\end{enumerate}
\begin{enumerate}[label=\Alph*)]
\item 1, 2.
\item faqat 3.
\item faqat 4.
\item 3, 4.
\item 1, 4.
\end{enumerate}


\end{enumerate}




\begin{center}
\begin{tabular}{|c|c|c|c|c|c|c|c|c|c|}
\hline
1a & 2e & 3e & 4a & 5c & 6d & 7e & 8d & 9a & 10c \\ \hline
11e & 12c & 13a & 14b & 15b & 16b & 17d & 18a & 19b & 20d \\ \hline
21e & 22a & 23d & 24b & 25d & 26c & 27d & 28c & 29c & 30d \\ \hline
31c & 32a & 33c & 34b & 35a & 36d & 37a & 38e & 39e & 40c \\ \hline
41c & 42a & 43e & 44b & 45a & 46d & 47e & 48d & 49d & \\ \hline
\end{tabular}
\end{center}